Historické šifrované dokumenty predstavujú jedinečný a neoceniteľný zdroj informácií o komunikácii, diplomacii a kultúre minulých storočí. Mnoho z týchto dokumentov však zostáva nerozlúštených, pretože tradičné metódy ich spracovania sú časovo náročné a vyžadujú špecializované znalosti. Rozvoj moderných technológií v oblasti spracovania obrazu, strojového učenia a webových aplikácií otvára nové možnosti pre systematickú digitalizáciu, anotáciu a analýzu týchto vzácnych materiálov.

Cieľom tohto tímového projektu bolo navrhnúť a implementovať kolaboratívnu webovú platformu, ktorá umožní efektívne spracovanie historických šifrovaných dokumentov. Platforma slúži ako nástroj pre historikov, kryptografov a výskumníkov, ktorí sa zaoberajú analýzou historických šifier. Systém kombinuje automatizované metódy spracovania obrazu a rozpoznávania znakov s možnosťou manuálnej anotácie a korekcie, čím vytvára polo-automatický pracovný postup optimalizovaný pre prácu s historickými dokumentmi.

Projekt vznikol v spolupráci s výskumnou komunitou zaoberajúcou sa historickou kryptografiou. Vedúci projektu, Ing. Stanislav Marochok, poskytol tímu prístup k datasetu historických šifrovacích kľúčov \textit{Transcendino} a definoval základné požiadavky na funkcionalitu systému vychádzajúce z reálnych potrieb výskumníkov v tejto oblasti.

Z technického hľadiska projekt zahŕňa vývoj viacerých vzájomne prepojených komponentov: frontendovej webovej aplikácie postavenej na knižnici React, backendového API servera implementovaného v jazyku Python s využitím frameworku FastAPI, databázového systému PostgreSQL a samostatného AI modulu pre spracovanie obrazu a detekciu štruktúry dokumentov. Celý systém je nasadený v kontajnerizovanom prostredí Docker a prevádzkovaný na vlastnom serverovom hardvéri, čo zabezpečuje jednoduchú správu, škálovateľnosť a prenositeľnosť riešenia.

Dokumentácia opisuje komplexné výsledky práce tímu počas oboch semestrov akademického roka 2025/2026 (predmety Tímový projekt 1 a Tímový projekt 2). V nasledujúcich kapitolách sú podrobne popísané požiadavky na systém, architektúra a implementácia jednotlivých komponentov vrátane šesťkrokového pipeline rozhrania pre spracovanie dokumentov, serverové prostredie a nasadenie aplikácie, experimentálne výsledky trénovania detekčných modelov, ako aj priebeh vývoja dokumentovaný prostredníctvom zápisníc z pravidelných stretnutí tímu.
